% 两个二元翻转核的复合。教学数据：第一核翻转概率 p=0.1，第二核翻转概率 q=0.2；
% 从 0 到 1 的两条路径 0->0->1 与 0->1->1 的概率为 0.9*0.2=0.18 与 0.1*0.8=0.08，
% 复合翻转概率 p(1-q)+(1-p)q=0.26。
\begin{tikzpicture}[line width=.6pt,font=\fontsize{8}{10}\selectfont,>=Stealth,
  st/.style={circle,draw=FigureBlue,fill=FigureBlueFill,minimum size=6mm,inner sep=0pt},
  mid/.style={circle,draw=FigureYellow,fill=FigureYellowFill,minimum size=6mm,inner sep=0pt},
  ot/.style={circle,draw=FigureGreen,fill=FigureGreenFill,minimum size=6mm,inner sep=0pt},
  lab/.style={fill=white,inner sep=1pt},
  note/.style={FigureMuted,font=\fontsize{7.5}{9}\selectfont}]
\node[st] (s0) at (0,0) {$0$};
\node[st] (s1) at (0,-16mm) {$1$};
\node[mid] (t0) at (32mm,0) {$0$};
\node[mid] (t1) at (32mm,-16mm) {$1$};
\node[ot] (u0) at (64mm,0) {$0$};
\node[ot] (u1) at (64mm,-16mm) {$1$};
\node[FigureBlue] at (0,8mm) {输入$s$};
\node[FigureYellow] at (32mm,8mm) {中间$t$};
\node[FigureGreen] at (64mm,8mm) {输出$u$};
\node[FigureInk] at (16mm,8mm) {核$K$};
\node[FigureInk] at (48mm,8mm) {核$L$};
% 第一核
\draw[->,FigureRed,line width=1.1pt] (s0) -- node[lab,above]{$1-p$} (t0);
\draw[->,FigureMuted!45,line width=.4pt] (s1) -- (t1);
\draw[->,FigureBlue,line width=1.1pt] (s0) -- node[lab,pos=.3]{$p$} (t1);
\draw[->,FigureMuted!45,line width=.4pt] (s1) -- (t0);
% 第二核
\draw[->,FigureMuted!45,line width=.4pt] (t0) -- (u0);
\draw[->,FigureBlue,line width=1.1pt] (t1) -- node[lab,below]{$1-q$} (u1);
\draw[->,FigureRed,line width=1.1pt] (t0) -- node[lab,pos=.3]{$q$} (u1);
\draw[->,FigureMuted!45,line width=.4pt] (t1) -- (u0);
% 路径汇总
\node[FigureRed,anchor=west] at (72mm,-2mm) {$0\to0\to1$：$(1-p)q=0.18$};
\node[FigureBlue,anchor=west] at (72mm,-9mm) {$0\to1\to1$：$p(1-q)=0.08$};
\node[FigureInk,anchor=west] at (72mm,-16mm) {$(KL)(0,\{1\})=0.26$};
\node[note,anchor=west] at (72mm,-23mm) {$p=0.1$，$q=0.2$};
\end{tikzpicture}
